\vfill

\tjzspis{Szczęście z~jelita} %tytuł do spisu treści
\tjztyt{Szczęście z~jelita} %tytuł odcinka 
Nie od dziś wiadomo, że ćwiczenia fizyczne budują formę i~poprawiają nastrój, zalecane są osobom narażonym na stres. Wiadomo także, że w~ciągłym ruchu są (lub chcą być) dzieci.  W~zeszłym roku naukowcy z~Szanghaju w~publikacji w~„Neuroscience” formułują hipotezę, że aktywność fizyczna wspomaga prawidłowy rozwój mózgu we wczesnym dzieciństwie, a~dzieje się tak dzięki serotoninie. Obserwując grupkę kilkulatków w~parku, zastanawiam się, ilu z~opiekunów zdaje sobie sprawę, że oto radosne brykanie jest tak naprawdę kształtowaniem funkcji poznawczych, rozwojem pamięci i~tworzeniem społeczno-emocjonalnej podstawy rozwijającego się człowieka. Jak maleńka cząsteczka zwana „hormonem szczęścia” bierze w~tym udział?

Autorzy pracy z~wieloletnich publikacji zebrali kilka najważniejszych informacji wyjściowych. Po pierwsze aktywność fizyczna u~dzieci korzystnie wpływa na ich rozwój: wzmacnia siłę mięśni, gęstość kości, poprawia wydolność krążeniowo-oddechową. Ale nie tylko.  Ruch służy  prawidłowemu rozwojowi kontroli emocji i~procesów poznawczych: myślenia, uczenia się, pamięci, i~tzw. „funkcji wykonawczych”, czyli tym, jak na podstawie dostępnych informacji i~wiedzy podejmować decyzje i~określone działania. Nic zatem dziwnego, że WHO w~2019 roku wydało rekomendację, w~której zaleca się, żeby dzieci między 3 a~4 rokiem życia minimum 3 godziny dziennie spędzały aktywnie fizycznie, w~tym godzinę intensywnie.
\eject
\begin{dwieszpalty}
Po drugie, od dawna badania wykazują, że aktywność fizyczna stymuluje przepływ sygnałów w~układzie nerwowym,  pobudza namnażanie się komórek nerwowych i~zwiększa produkcję neuroprzekaźników.

Po trzecie,  autorzy cytują prace wykazujące, że aktywność fizyczna podnosi poziom serotoniny. Przynajmniej u~zwierząt: wprowadza się im do mózgu cienkie rurki, umożliwiające pobieranie próbek do analizy, a~następnie wkłada do koła treningowego kręcącego się z~dużą prędkością przez dwie godziny. Bez względu na to, co myślimy o~losie szczurów laboratoryjnych, trzeba przyznać, że wyniki tych badań wskazują, iż poziom serotoniny w~mózgach ćwiczących zwierząt rośnie. Problem  polega na tym, że trudno bada się poziom czegokolwiek w~ludzkim mózgu, a~szczególnie jeśli jest to mózg małego dziecka.

Dlaczego nie pobrać po prostu krwi? Tu się sprawy komplikują.

Serotonina to niewielka cząsteczka, która powstaje z~jednego z~aminokwasów (cegiełek budulcowych białek),  tryptofanu. Ludzkie komórki nie potrafią syntetyzować tryptofanu, pochodzi on zatem z~trawionego jedzenia (głównie mięsa, nabiału, nasion soi i~dyni).

Serotonina powstaje z~tryptofanu w~dwóch reakcjach, z~których pierwsza prowadzona jest przez  enzym TPH. Ponad 90\% serotoniny w~ludzkim ciele powstaje w~tzw. komórkach chromochłonnych wyściółki żołądka, jelita cienkiego i~grubego, dzięki aktywności enzymu TPH1. Serotonina trafia do krwi obwodowej i~jest magazynowana w~płytkach krwi. Reguluje tu krzepnięcie krwi,  ciśnienie tętnicze (dzięki wpływowi na kurczenie się naczyń krwionośnych), pobudza prawidłową pracę żołądka i~jelit. Ale w~mózgu cząsteczki serotoniny pełnią inne funkcje. Jest neuroprzekaźnikiem – wpływa zatem na przekazywanie sygnałów między neuronami, reguluje działanie wielu obszarów mózgu, głównie tych odpowiadających za regulacje emocji, snu, apetytu, wpływa na pamięć, koncentrację uwagi, zachowania seksualne. Najważniejsze jednak, że jest ona niezbędna do prawidłowego rozwoju mózgu, szczególnie w~okresie prenatalnym i~kilku pierwszych latach życia. Już po kilku tygodniach od zapłodnienia rozwój układu nerwowego, a~szczególnie mózgu, wskakuje na najwyższy możliwy poziom. Według niektórych szacunków w~tym czasie powstaje kilkadziesiąt tysięcy komórek nerwowych na minutę. Mózg 5-latka waży prawie tyle samo co osoby dorosłej. Nie świadczy to jednak o~poziomie jego rozwoju: nie chodzi tylko o~liczbę komórek, ale przede wszystkim o~komunikację między nimi. Po czasie szaleńczego wzrostu liczby komórek nerwowych  następuje ich dojrzewanie i~gwałtowne zwiększanie się liczby, czyli połączeń między nimi. Szacuje się, że w~pierwszych latach życia może ich przybywać nawet milion na sekundę. Po tym czasie następuje proces nazwany po angielsku „pruning” –  przycinanie: liczba połączeń między neuronami zmniejsza się, pozostają tylko te, które były najbardziej używane. W~procesach tych serotonina jest niezbędna do powstawania nowych komórek nerwowych, do przemieszczania się niedojrzałych neuronów w~odpowiednie miejsca, do różnicowania się komórek na różne typy i~ich dojrzewania,  w~końcu stymuluje powstawanie synaps i~decyduje o~tzw. plastyczności synaptycznej – moduluje ich działanie.

Serotonina mózgowa nie pochodzi z~puli cząsteczek krwi obwodowej, ponieważ nie przechodzi przez barierę krew-mózg.  Na szczęście barierę tę przechodzi tryptofan, i~z~niego powstaje serotonina w~tzw. jądrach szwu w~rejonie pnia mózgu.
\end{dwieszpalty}

W~finale dochodzenia chińskich naukowców aktywność fizyczna wpływa nie tyle na samą serotoninę, ale na  obecność we krwi jej wolnego prekursora, tryptofanu.  Autorzy postulują dwa możliwe mechanizmy. Po pierwsze, tryptofan we krwi wiąże się z~białkiem, albuminą, z~której jest uwalniany w~razie potrzeby. Albumina transportuje także wolne kwasy tłuszczowe (tzw. NEFA), których pojawia się dużo przy długotrwałym lub intensywnym wysiłku fizycznym. NEFA i~tryptofan konkurują ze sobą. Kiedy cząsteczek NEFA jest więcej, wypierają z~albuminy tryptofan, a~ten, uwolniony, przenika do mózgu.

Po drugie, aktywność fizyczna wpływa na florę bakteryjną w~jelitach. Niektóre gatunki mikroorganizmów syntetyzują i~wydzielają poza komórki tryptofan (i~także serotoninę). Gatunki te szczególnie dobrze rozwijają się w~jelitach ludzi aktywnych fizycznie, także dzieci, szczególnie tych, które mają kontakt z~naturą, gdzie mogą się ubrudzić i~jednocześnie zmęczyć. W~ten sposób, bez względu na dietę, źródłem tryptofanu wędrującego do mózgu są bakterie jelitowe.

Wciąż pokutuje podział sprzed wielu lat na dzieci, które się uczą i~na te, które biegają. Wygląda jednak na to, że to ci, którzy  biegają, skaczą, wspinają się, dyndają, jeżdżą, kręcą się i~fikają koziołki będą się uczyć najlepiej. Czyżbyśmy zapomnieli, że jesteśmy skąpo owłosioną małpą, która niedawno zeszła z~drzewa?

 
\medskip
\rightline{\large\it Marta FIKUS-KRYŃSKA}

\eject
